\section{The Tengwar}
Phi can be written using the standard Latin alphabet, but is usually written
using J. R. R. Tolkien's
Tengwar\footnote{\url{https://en.wikipedia.org/wiki/Tengwar}} script. The Tengwar
was first described in Appendix E of the third book in Tolkein's trilogy,
\emph{The Lord of the Rings}. There are numerous sources of information
available concerning the Tengwar, so the history and details of the script will
not be explored in this document.

The look and feel of the Tengwar's characters (\emph{tengwa}) and its
diacritical marks (\emph{tehta}) match the ideals and philosophies of this
language. Phi has a small set of consonants and vowels so it is relatively
straightforward to map its graphemes to Tengwar characters. The Unicode font
\emph{Tengwar Telcontar}\footnote{\url{http://freetengwar.sourceforge.net/%
tengtelc.html}} is used in this manual to render Phi text in Tengwar script. A
custom keyboard layout\footnote{\url{http://freetengwar.sourceforge.net/%
keylayouts.html}} was employed to facilitate typing Tengwar characters.

Tables \ref{tab:tengwar-consonants} and \ref{tab:tengwar-vowels} outline the
mapping between the graphemes used in Phi and the corresponding Tengwar
characters, as well as the keypresses necessary to produce the characters using
the font and the keyboard layout mentioned previously. An application or utility
that supports Graphite fonts\footnote{\url{https://scripts.sil.org/cms/%
scripts/page.php?site\_id=projects\&item\_id=graphite\_fonts}} is required for
the font to render correctly. A link to details of this requirement can be found
in the second footnote below.

Phi uses the Tengwar script in its own mode. For clarity this will be referred
to as the Phi mode. A couple of departures from other Tengwar modes are in the
application of the
\emph{r}~rule\footnote{\url{https://www.wikihow.com/Write-Tengwar\#step-id-16}},
and the general suggestion that \emph{tehta} not be placed below \emph{tengwa}.
The first is a minor departure; the makeup of Phi words is very regular so the
\emph{r} rule is modified for this situation. See below for specifics.

The second is a larger departure. Words in Phi often contain two consecutive
vowels. The first \emph{tehta} is placed above the consonant that precedes the
vowel, as in the Quenya mode. But to avoid the use of a vowel carrier, the
second \emph{tehta} is placed below the same consonant. The letters would then
be read as consonant, upper vowel, then lower vowel.
\clearpage
This \emph{tehtar} strategy and the use of a single Tengwar character for
consonant doubles such as \dscopy ‹th›~\normalfont or \dscopy ‹sh›~\normalfont
cause Phi script to be very compact and efficient. To illustrate, the word for
\enquote{love} in Phi is \emph{lothea}. In the Tengwar script it looks like
this:

\bigskip
\leftskip=0.66in
\ttlarge 
\normalfont
\bigskip

\leftskip=0in
This six-letter word in latin characters requires only two \emph{tengwa} to
represent it. Many words in Phi are two to three letters long. This occurs
mainly in parts of speech such as particles, prepositions, or pronouns. These
words are almost always a single \emph{tengwa} with diacriticals in Tengwar
script. All of this results in a very concise writing system.

This is the phrase \enquote{I love you} in its simplest vernacular form in Phi:

\bigskip
\leftskip=0.66in
\ttlarge 
\normalfont
\bigskip

\leftskip=0in
That is a beautiful sentiment rendered in a beautiful script. For more formal
communication in Phi, particles are used to mark parts of speech in order to
reduce ambiguity. As well, word markers are often used to distinguish
multisyllabic words for clarity. The Phi Tengwar punctuation is outlined in
Table~\ref{tab:tengwar-punctuation}. As an example, the phrase \enquote{I love
you} would be written formally as:

\bigskip
\leftskip=0.66in
\ttlarge ⸱⸱⸱⸱⸱
\normalfont
\bigskip

\leftskip=0in
In the Phi consonant mapping illustrated in Table~\ref{tab:tengwar-consonants}
there are two \emph{tengwa} that represent the \dscopy [s]~\normalfont sound:
\enquote{\ttcopy \normalfont} and \enquote{\ttcopy \normalfont.} Which of the
two is used in a word is at the writer's discretion. Most often it depends on
which diacritical marks, if any, are attached to the character. If the \dscopy
‹s›~\normalfont \emph{tengwa} looks uncomfortable in one orientation, then the
alternate might be more effective.

As well, there are two \emph{tengwa} to represent \dscopy /r/\normalfont:
\enquote{\ttcopy \normalfont} and \enquote{\ttcopy \normalfont.} In this case,
however, there is a specific usage for both forms. In Phi the initial \dscopy
/r/~\normalfont in a word is always a trilled \dscopy [r]~\normalfont and is
represented by \enquote{\ttcopy \normalfont}. Within a word \dscopy /r/
~\normalfont is always a tapped \dscopy [ɾ]~\normalfont and is represented by
\enquote{\ttcopy \normalfont}.

\leftskip=0in
\clearpage
\begin{table}[H]
  \caption{\label{tab:tengwar-consonants}Consonant grapheme to tengwa mapping}
  \renewcommand\arraystretch{2.5}
  \setlength\extrarowheight{-3pt}
  \begin{tabular}{|c|c|c|}
    \hline
    \rowcolor{greyhead}
    Grapheme & Tengwa & Keypress \\
    \hline
    \dsil ‹p› & \ttsmall  & p \\
    \hline
    \rowcolor{greyrow}
    \dsil ‹t› & \ttsmall  & t \\
    \hline
    \dsil ‹m› & \ttsmall  & m \\
    \hline
    \rowcolor{greyrow}
    \dsil ‹n› & \ttsmall  & n \\
    \hline
    \dsil ‹r› & \ttsmall  & r \\
    \hline
    \rowcolor{greyrow}
    \dsil ‹r› & \ttsmall  & shift-r \\
    \hline
    \dsil ‹ph› & \ttsmall  & shift-p \\
    \hline
    \rowcolor{greyrow}
    \dsil ‹wh› & \ttsmall  & shift-k \\
    \hline
    \dsil ‹th› & \ttsmall  & shift-t \\
    \hline
    \rowcolor{greyrow}
    \dsil ‹s› & \ttsmall  & s \\
    \hline
    \dsil ‹s› & \ttsmall  & shift-s \\
    \hline
    \rowcolor{greyrow}
    \dsil ‹sh› & \ttsmall  & shift-c \\
    \hline
    \dsil ‹h› & \ttsmall  & h \\
    \hline
    \rowcolor{greyrow}
    \dsil ‹w› & \ttsmall  & w \\
    \hline
    \dsil ‹l› & \ttsmall  & l \\
    \hline
  \end{tabular}
\end{table}
\begin{table}[H]
  \caption{\label{tab:tengwar-vowels}Vowel grapheme to tehta mapping}
  \renewcommand\arraystretch{2.5}
  \setlength\extrarowheight{-3pt}
  \begin{tabular}{|c|c|c|}
    \hline
    \rowcolor{greyhead}
    Grapheme &
    \hspace{1em} Tehta \hspace{1em} &
    \hspace{1em} Keypress \hspace{1em} \\
    \hline
    \dsil ‹i› & \ttsmall   & i / option-i\footnotemark \\
    \hline
    \rowcolor{greyrow}
    \dsil ‹u› & \ttsmall   & u / option-u \\
    \hline
    \dsil ‹e› & \ttsmall   & e / option-e \\
    \hline
    \rowcolor{greyrow}
    \dsil ‹o› & \ttsmall   & o / option-o \\
    \hline
    \dsil ‹a› & \ttsmall   & a / option-a \\
    \hline
  \end{tabular}
\end{table}
\footnotetext{Keypresses for the lower diacriticals are given for macOS. Use
\emph{ctrl-alt} on Windows instead of \emph{option}.}
\begin{table}[H]
  \caption{\label{tab:tengwar-punctuation}Punctuation to tengwa mapping}
  \renewcommand\arraystretch{2.5}
  \setlength\extrarowheight{-3pt}
  \begin{tabular}{|c|c|c|}
    \hline
    \rowcolor{greyhead}
    Punctuation & Tengwar & Keypress \\
    \hline
    space & \ttsmall ⸱ & , \\
    \hline
    \rowcolor{greyrow}
    period & \ttsmall : & shift-, \\
    \hline
  \end{tabular}
\end{table}
